\documentclass{article}
\usepackage{fontspec}
\setmainfont{Microsoft YaHei}
\XeTeXlinebreaklocale "zh"
\XeTeXlinebreakskip = 0pt plus 1pt
\usepackage[margin=2.2cm]{geometry}
\usepackage{booktabs}
\usepackage{enumitem}
\setlist{nosep}
\title{阶段报告：环境修复、GPU 实验与真实数据接入}
\author{}
\date{2026-09-24}
\begin{document}
\maketitle

\section{概要}
本阶段完成四项工作。其一，修复计算环境：升级 torch 至 2.13.0+cu126 并使 torchvision 与其配对，期间发现并替换了一枚 CPU 构建的 torchvision（该构建不含 CUDA 算子内核，无法调用显卡），修复后通过 19 项核验，升级前后数值逐位一致。其二，搭建真 GPU 渲染器并重跑五组跟踪实验，三条判据全部通过：末端误差 0.8--1.0 毫米，每帧 0.2--0.8 秒。其三，随机种子复核通过并完整存档。其四，接入 Apollo 真实数据集：还原其考卷结构、取得官方对照成绩、绕过缺失的模型文件完成真值组装，并在真实照片上完成打分器判别力抽查。

\section{一、计算环境修复与核验}
环境中原存两处 2026 年 7 月遗留问题：setuptools 版本超出 torch 的声明区间；torchvision 0.28.0 要求 torch==2.13.0 而实装为 2.12.0。核验中进一步发现 torchvision 实为 CPU 构建（版本号后缀 +cpu，其 CUDA 算子无注册内核，实调即报错）——这正是历来防护条款所要杜绝的情形。处置：torch 升至 2.13.0+cu126、torchvision 替换为 0.28.0+cu126（CUDA 构建）、setuptools 降至 81.0.0（满足 torch 声明的 $\geq$77.0.3）。

升级前后以固定输入做数值对照：矩阵乘法结果最大相对差 $0.00\times10^{0}$（逐位一致），视觉模型同输入输出余弦相似度 1.000000——升级对数值零扰动。核验矩阵 19 项判据全部通过，涵盖版本配对、显卡适配（RTX 4060 / 计算能力 8.9 / 运行时 CUDA 12.6 与驱动 581.80 向后兼容）、数值血统对照、真实负载端到端（视觉模型前向、torchvision CUDA 算子实调、依赖冒烟）。全程 numpy 2.4.4 未变动；回滚轮子与 SHA256 指纹已存档备命。

\section{二、GPU 渲染与模拟实验}
上午核验确认 pybullet 的 GPU 渲染路径不可用（EGL 插件缺失；硬件 OpenGL 静默回退为 CPU 软渲染，与软渲染逐像素一致）。遂自建渲染器：以 moderngl 做显卡光栅化绘制机械臂，运动学由 pybullet 提供，网格经 trimesh 读取。实测 128$^2$ 分辨率下 0.5 毫秒/帧，为 CPU 软渲染（28.2 毫秒）的 56 倍；224$^2$ 下快 138 倍。三判据（真显卡出图、同姿态画面质心对齐 4.3 像素、快于 CPU 基线）全部通过。

随后将该渲染器接入粒子滤波实验，重跑 300 帧五组对照。数据轨迹、预埋偏差、随机种子 20260924 与算法参数均与 CPU 版逐位同源，唯一变量是渲染器。结果见表~\ref{tab:gpu}。

\begin{table}[h]\centering
\caption{GPU 渲染版实验结果（300 帧，末端位置误差，毫米）}
\label{tab:gpu}
\begin{tabular}{lccccc}
\toprule
方法 & 中位 & 均值 & 95分位 & 每帧耗时 & 峰值显存 \\
\midrule
死推算（不看图） & 24.1 & 24.4 & 28.1 & --- & --- \\
粒子滤波-像素-带提示 & \textbf{0.8} & 2.1 & 6.9 & 194ms & --- \\
粒子滤波-像素-随机走 & 11.8 & 12.6 & 21.7 & 203ms & --- \\
粒子滤波-DINOv2-带提示 & \textbf{1.0} & 2.9 & 15.0 & 822ms & 0.84GB \\
粒子滤波-DINOv2-随机走 & 14.3 & 14.6 & 23.7 & 844ms & 0.84GB \\
\bottomrule
\end{tabular}
\end{table}

与 CPU 版对照（上周结果：12.8 / 18.0 / 13.0 / 76.8 毫米），误差大幅下降，三条判据由"仅先验价值一条通过"变为全部通过：主判据（误差 $\leq$10 毫米且优于基线一半）以 0.8 与 1.0 毫米通过（优于基线 96\%--97\%）；先验价值（带提示优于随机走三成）两组均好 93\%；笔记本可行（每帧 $\leq$3 秒）以 822 毫秒通过。因果解释：渲染画面的边缘更干净，打分器分辨力随之变锐，偏差辨识精度由 1--2 度提升至 0.1 度量级；上周 12.8 毫米的残留误差，根源在画面锐度而非算法。

\section{三、随机种子复核与存档}
以同一随机种子 20260924 重跑实验前 100 帧，五个检查点的误差与有效粒子数和存档逐点一致（14.0 / 9.5 / 16.9 / 16.8 / 24.9 毫米；有效粒子数 52 / 52 / 44 / 62 / 60）。全部随机源已穷尽清查：仅 run\_pf.py 使用随机数，五组配置各自持有同值种子（配对设计）；数据生成脚本零随机。种子表、消耗顺序、五个基准文件的 SHA256 指纹与复核记录存档于 \_mvp/SEEDS.md。

\section{四、Apollo 真实数据集}
该数据集为马普所 2016 年 hand\_tracking\_dataset（数据分区代号 Apollo）：固定基座双 KUKA LWR IV 机械臂加 Barrett 三指灵巧手，头部装主动人形颈与 RGB-D 相机，VICON 动捕为标准答案。考卷含 12 段真实序列（7 段 60 秒加 5 段 30 秒正弦运动），另附两套改写关节读数的版本——恒定偏差每关节 8.6 度、方波偏差 5 度——图像与真值保持原样。考核指标为手对相机六自由度位姿的平移误差与角度误差。

\textbf{官方对照成绩}（其 RA-L 2016 论文《Probabilistic Articulated Real-Time Tracking for Robot Manipulation》在恒定偏差卷上的中位误差，表~\ref{tab:ra}）：融合路线 0.62 厘米，纯编码器 19.7 厘米，纯视觉无先验 47.8 厘米——运动先验的价值在真实数据上有官方证据（好约 99\%），本课题的对照基线由此从概念变为具体数字。

\begin{table}[h]\centering
\caption{官方基准（synthetic\_constant\_bias，中位误差）}
\label{tab:ra}
\begin{tabular}{lcc}
\toprule
路线 & 平移误差 & 角度误差 \\
\midrule
纯编码器（ROTARY） & 19.7 cm & 27.0$^\circ$ \\
纯视觉变体（CAMERA） & 14.0 cm & 26.4$^\circ$ \\
融合（FUSION） & \textbf{0.62 cm} & \textbf{2.3$^\circ$} \\
纯视觉无先验（VISUAL） & 47.8 cm & 79.8$^\circ$ \\
\bottomrule
\end{tabular}
\end{table}

\textbf{模型缺失的绕行}：项目宣称提供 URDF 与开源代码，实查为烂尾（仓库仅数据压缩包与文档，代码检索三路无果）。改用数据自足路线：消息流中的整机 TF 骨骼为实测值（22 对刚性坐标系，双观测散布 0.000 毫米），手部骨架按实测偏移拼装、厚度取近似；真值手对相机位姿由动捕双刚体加一次标定组装，标定双观测一致性 0.10 毫米 / 0.03 度。以该口径在真实序列上测得纯运动学位姿误差 0.68 厘米 / 1.12 度（即无偏差时的误差地板）。

\textbf{真实照片判别力抽查}（20 帧 $\times$ 37 候选位姿，判据预先定死并按扰动类型分轴判定）：平移维两代打分器均通过（8 厘米大扰动逐帧命中率 85.8\%（凸包 IoU）与 80.8\%（深度残差），扰动量与分数单调下降）；旋转维不通过（大扰动命中率 45.8\% 与 54.2\%，接近随机）——但平均分曲线两代均单调恶化，二代深度残差的旋转信号幅度为一代的 20 倍。结论分层：真实照片上平移判别成立；旋转判别受限于手部模型的厚度近似（缺指形与掌形不对称），非打分器或照片之过。补精细手部模型（Barrett Hand 有公开规格）即可继续检验旋转维。

\section{五、判据总账}
\begin{table}[h]\centering
\begin{tabular}{lcc}
\toprule
判据组 & 结果 & 备注 \\
\midrule
模拟实验（三条判据） & 3/3 通过 & 上周 1/3，本周边界条件未变 \\
环境核验（19 项） & 19/19 通过 & 含数值血统逐位对照 \\
抽查-平移维（两条） & 2/2 通过 & 两代打分器均过 \\
抽查-旋转维（两条） & 0/2 通过 & 病根=模型保真度，记录在案 \\
\bottomrule
\end{tabular}
\end{table}

\section{六、结论与下一步}
本阶段证明：视觉修正机制不仅在自渲染画面下成立，其核心判别力（平移维）在真实照片上同样成立；配合 GPU 渲染后，模拟实验的精度与实时性判据同时达标，官方基准表给出了真实数据上的量化目标（融合路线 0.62 厘米）。

下一步三选一：（1）以公开规格制作精细手部模型，补齐旋转维判别力检验；（2）在恒定偏差考卷上运行完整跟踪流程，与官方成绩对照；（3）整理中期报告正文。

\end{document}
